Vision-language models (VLMs) show strong multimodal capabilities but still struggle with fine-grained vision-language reasoning.
We find that long chain-of-thought (CoT) reasoning exposes diverse failure modes, including perception, reasoning, knowledge, and hallucination errors, which can compound across intermediate steps.
However, most existing vision-language data used for reinforcement learning with verifiable rewards (RLVR) does not involve complex reasoning chains that rely on visual evidence throughout, leaving these weaknesses largely unexposed.
We therefore propose \textbf{HopChain}, a scalable framework for synthesizing \emph{multi-hop vision-language reasoning} data for RLVR training of VLMs.
Each synthesized multi-hop query forms a logically dependent chain of instance-grounded hops, where earlier hops establish the instances, sets, or conditions needed for later hops, while the final answer remains a specific, unambiguous number suitable for verifiable rewards.
We train Qwen3.5-35B-A3B and Qwen3.5-397B-A17B under two RLVR settings: the original data alone, and the original data plus HopChain's multi-hop data, and compare them across 24 benchmarks spanning STEM and Puzzle, General VQA, Text Recognition and Document Understanding, and Video Understanding.
Although this multi-hop data is not synthesized for any specific benchmark, it improves 20 of 24 benchmarks on both models, indicating broad and generalizable gains.
Consistently, replacing full chained queries with half-multi-hop or single-hop variants reduces the average score across five representative benchmarks from 70.4 to 66.7 and 64.3, respectively.
Notably, multi-hop gains peak in long-CoT vision-language reasoning, exceeding 50 points in the ultra-long-CoT regime.
These experiments establish HopChain as an effective, scalable framework for synthesizing multi-hop data that improves generalizable vision-language reasoning.
